% oss-plantilla: oss-2026.11
% Preámbulo corporativo fijo. No editar en sd-NN.tex: exportar_latex.py verificar lo compara byte a byte.
% Para cambiar el formato, editar esta plantilla u oss.sty y subir la versión de la primera línea.
\PassOptionsToPackage{unicode}{hyperref}
\PassOptionsToPackage{hyphens}{url}
\documentclass[11pt]{article}
\usepackage{amsmath,amssymb}
\usepackage{iftex}
\ifPDFTeX
  % Prism y otros editores en la nube: pdfLaTeX con Helvetica como familia por defecto.
  \usepackage[T1]{fontenc}
  \usepackage[utf8]{inputenc}
  \usepackage{textcomp}
  \usepackage[scaled=0.92]{helvet}
  \renewcommand{\familydefault}{\sfdefault}
  \DeclareUnicodeCharacter{2264}{\ensuremath{\leq}}
  \DeclareUnicodeCharacter{2265}{\ensuremath{\geq}}
  \DeclareUnicodeCharacter{2260}{\ensuremath{\neq}}
  \DeclareUnicodeCharacter{2248}{\ensuremath{\approx}}
  \DeclareUnicodeCharacter{2190}{\ensuremath{\leftarrow}}
  \DeclareUnicodeCharacter{2192}{\ensuremath{\rightarrow}}
  \DeclareUnicodeCharacter{25CF}{\ensuremath{\bullet}}
  \DeclareUnicodeCharacter{2713}{\ensuremath{\checkmark}}
\else
  % XeLaTeX/LuaLaTeX (compilación oficial local): DejaVu Sans.
  \usepackage{fontspec}
  \defaultfontfeatures{Ligatures=TeX,Scale=MatchLowercase}
  \setmainfont{DejaVu Sans}
  \setsansfont{DejaVu Sans}
  \setmonofont{DejaVu Sans Mono}
\fi
\usepackage[letterpaper,margin=20mm]{geometry}
\usepackage{microtype}
\usepackage{parskip}
\usepackage{xcolor}
\usepackage{longtable,booktabs,array}
\usepackage{calc}
\usepackage{etoolbox}
\makeatletter
\patchcmd\longtable{\par}{\if@noskipsec\mbox{}\fi\par}{}{}
\makeatother
\usepackage{footnotehyper}
\makesavenoteenv{longtable}
\usepackage{graphicx}
\makeatletter
\newsavebox\pandoc@box
\newcommand*\pandocbounded[1]{%
  \sbox\pandoc@box{#1}%
  \Gscale@div\@tempa{\textheight}{\dimexpr\ht\pandoc@box+\dp\pandoc@box\relax}%
  \Gscale@div\@tempb{\linewidth}{\wd\pandoc@box}%
  \ifdim\@tempb\p@<\@tempa\p@\let\@tempa\@tempb\fi
  \ifdim\@tempa\p@<\p@\scalebox{\@tempa}{\usebox\pandoc@box}%
  \else\usebox{\pandoc@box}%
  \fi}
\def\maxwidth{\ifdim\Gin@nat@width>\linewidth\linewidth\else\Gin@nat@width\fi}
\def\maxheight{\ifdim\Gin@nat@height>\textheight\textheight\else\Gin@nat@height\fi}
\makeatother
\setkeys{Gin}{width=\maxwidth,height=\maxheight,keepaspectratio}
\makeatletter
\def\fps@figure{htbp}
\makeatother
\usepackage[normalem]{ulem}
\providecommand{\st}[1]{\sout{#1}}
\usepackage{fancyvrb}
\setlength{\emergencystretch}{3em}
\providecommand{\tightlist}{%
  \setlength{\itemsep}{0pt}\setlength{\parskip}{0pt}}
\providecommand{\passthrough}[1]{#1}
\setcounter{secnumdepth}{-\maxdimen}
\ifPDFTeX
  \usepackage[spanish]{babel}
\else
  \usepackage[bidi=default]{babel}
  \babelprovide[main,import]{spanish}
\fi
\let\LanguageShortHands\languageshorthands
\def\languageshorthands#1{}
\usepackage{bookmark}
\usepackage{xurl}
\urlstyle{same}
\hypersetup{pdflang={es-CL},pdfcreator={Only Simple Solutions - exportar_latex.py}}
\usepackage{oss}
\begin{document}

\ossCover{Formulario T-11}{Subdocumento 4}
\tableofcontents
\clearpage

\renewcommand{\tablename}{Tabla}
\renewcommand{\thetable}{T-11.\arabic{table}}

\section{Formulario T-11 · Especificaciones técnicas ofertadas}\label{formulario-t-11}

Conforme al Formulario T-7 y al Formulario T-21 de las Bases Administrativas, este formulario detalla los componentes de infraestructura, plataforma, licenciamiento y hardware de la solución ofertada a Multitiendas Ancoa S.A., con su ubicación y su justificación. Las secciones 4.1.1, 4.2, 4.2.1 y 4.3 del Subdocumento 4 contienen el resumen y el análisis. Este formulario reúne el listado completo, incluidos los componentes que el Subdocumento 4 menciona en el texto y en las figuras sin incorporarlos a la Tabla 4.6.

El formulario usa los códigos de sitio de la Tabla 4.8 del Subdocumento 4. SIT-01 agrupa las 22 tiendas, CD-01 es el centro de distribución principal, CD-02 es el centro de distribución de Concepción, DC-01 es el centro de datos de casa matriz, DC-02 es su sala de respaldo, CLD-01 es la región Azure Chile Central, CLD-02 es la región Azure Brazil South y EXT-01 reúne a las contrapartes externas. Cada componente se atribuye a un ámbito, que puede ser Retail, Emisor o plataforma compartida, porque la separación entre la tienda y la Filial emisora alcanza a la infraestructura. La columna Origen indica el documento de la propuesta que selecciona el componente. Cuando el valor es T-11, este formulario precisa el producto o el nivel de una capacidad que los demás documentos dejaron remitida a él.

La Tabla~\mbox{\ref{tab:t11-1}} indica qué contiene cada apartado y con qué sección de la propuesta se corresponde.

\begin{longtable}[]{@{}
  >{\raggedright\arraybackslash}p{(\columnwidth - 4\tabcolsep) * \real{0.3000}}
  >{\raggedright\arraybackslash}p{(\columnwidth - 4\tabcolsep) * \real{0.4500}}
  >{\raggedright\arraybackslash}p{(\columnwidth - 4\tabcolsep) * \real{0.2500}}@{}}
\caption{Correspondencia entre el formulario y la propuesta.}\label{tab:t11-1}\\
\toprule\noalign{}
\textbf{Apartado} & \textbf{Contenido} & \textbf{Sección de la propuesta} \\
\midrule\noalign{}
\endfirsthead
\toprule\noalign{}
\textbf{Apartado} & \textbf{Contenido} & \textbf{Sección de la propuesta} \\
\midrule\noalign{}
\endhead
\bottomrule\noalign{}
\multicolumn{3}{@{}p{\columnwidth}@{}}{\rule{0pt}{3.2ex}\textbf{Fuente:} Elaboración propia.} \\
\endlastfoot
T-11.1 Servicios de nube contratados & Servicios de Azure por región, nivel, cantidad, ámbito y justificación. & SD-04, 4.1.1 y 4.2 \\ \addlinespace
T-11.2 Plataforma y software de la solución & Lenguajes, marcos, motores y software propio, con versión. & SD-04, 4.1.1 \\ \addlinespace
T-11.3 Seguridad, entrega y operación & Cadena de entrega, análisis de seguridad, monitoreo y respuesta. & SD-04, 4.2; SD-06, 6.2; Formulario T-10 \\ \addlinespace
T-11.4 Hardware que provee el proponente & Nodos locales de tienda y sus repuestos. & SD-04, 4.2.1 \\ \addlinespace
T-11.5 Hardware y obras que adquiere el Cliente & Especificación de red de sitio, gabinete, energía, periféricos y terminales. & SD-04, 4.2.1 \\ \addlinespace
T-11.6 Conectividad & Enlaces por sitio, conexión privada y respaldo. & SD-04, 4.2 \\ \addlinespace
T-11.7 Centros de datos & Centro primario, centro secundario e instalaciones del Cliente. & SD-04, 4.3 \\ \addlinespace
T-11.8 Plataformas conservadas & Plataformas del Cliente que la solución integra y no oferta. & SD-04, 4.1 y 4.2.1 \\ \addlinespace
T-11.9 Componentes de las innovaciones & Equipos y plataformas de las cinco innovaciones. & SD-13 \\ \addlinespace
T-11.10 Licenciamiento & Modalidad, licencia, titular y vigencia por producto. & SD-04, 4.2.1 \\ \addlinespace
T-11.11 Ciclo de vida y soporte & Versión ofertada, fin de soporte publicado y plan de actualización. & SD-04, 4.1.1; SD-08 \\ \addlinespace
T-11.12 Portabilidad & Componentes portables y componentes propios del proveedor. & SD-04, 4.2; SD-08 \\ \addlinespace
T-11.13 Comprobaciones previas a la construcción & Verificaciones que fijan cantidades, modelos y niveles, con su hito. & SD-04, 4.2; SD-07 \\
\end{longtable}

Este formulario no contiene precios, tarifas ni montos que permitan inferir la Oferta Económica. Las cantidades que declara son las que alimentan esa oferta.

\subsection{T-11.1 Servicios de nube contratados}\label{t-11-1}

Microsoft Azure es el proveedor de nube (ADR-07). Los servicios se contratan bajo una suscripción del proyecto por ambiente y por ámbito, con Chile Central como región primaria y Brazil South como región de recuperación, ambas con tres zonas de disponibilidad. La Tabla~\mbox{\ref{tab:t11-2}} lista cada servicio con su nivel, su cantidad en el ambiente de Producción y la razón de su elección. Los cinco ambientes (Desarrollo, QA, Preproducción, Producción y Recuperación ante Desastres) usan las mismas clases de componentes. Preproducción y Recuperación tienen capacidades ajustadas a su finalidad y verificadas en la prueba de carga.

{\footnotesize
\begin{longtable}[]{@{}
  >{\raggedright\arraybackslash}p{(\columnwidth - 12\tabcolsep) * \real{0.0400}}
  >{\raggedright\arraybackslash}p{(\columnwidth - 12\tabcolsep) * \real{0.1200}}
  >{\raggedright\arraybackslash}p{(\columnwidth - 12\tabcolsep) * \real{0.1700}}
  >{\raggedright\arraybackslash}p{(\columnwidth - 12\tabcolsep) * \real{0.1600}}
  >{\raggedright\arraybackslash}p{(\columnwidth - 12\tabcolsep) * \real{0.1500}}
  >{\raggedright\arraybackslash}p{(\columnwidth - 12\tabcolsep) * \real{0.2400}}
  >{\raggedright\arraybackslash}p{(\columnwidth - 12\tabcolsep) * \real{0.1200}}@{}}
\caption{Servicios de Azure contratados.}\label{tab:t11-2}\\
\toprule\noalign{}
\textbf{N.°} & \textbf{Componente} & \textbf{Producto y nivel} & \textbf{Cantidad en Producción} & \textbf{Ubicación y ámbito} & \textbf{Justificación} & \textbf{Origen} \\
\midrule\noalign{}
\endfirsthead
\toprule\noalign{}
\textbf{N.°} & \textbf{Componente} & \textbf{Producto y nivel} & \textbf{Cantidad en Producción} & \textbf{Ubicación y ámbito} & \textbf{Justificación} & \textbf{Origen} \\
\midrule\noalign{}
\endhead
\bottomrule\noalign{}
\multicolumn{7}{@{}p{\columnwidth}@{}}{\rule{0pt}{3.2ex}\textbf{Fuente:} Elaboración propia a partir del Subdocumento 4 (Tablas 4.6, 4.7 y 4.10, y secciones 4.2 y 4.3) y del Formulario T-10.} \\
\endlastfoot
1.1 & Suscripciones & Suscripción de Azure del proyecto. & Una por ambiente y por ámbito. & CLD-01 y CLD-02. Ámbito Retail, Emisor y compartida. & Una credencial o una política no sirve en otro ambiente ni en el otro ámbito. & SD-04, 4.2 \\ \addlinespace
1.2 & Microservicios propios y adaptadores & Azure Kubernetes Service, nivel Standard, con nodos repartidos en tres zonas. & Un clúster por ámbito; nodos según la memoria de capacidad. & CLD-01; se crea en CLD-02 con la misma definición de infraestructura. Ámbito Retail y Emisor. & Servicio administrado de contenedores con zonas y escalado horizontal; evita operar el plano de control. & SD-04, Tabla 4.6; nivel T-11 \\ \addlinespace
1.3 & Intermediario de eventos & Azure Event Hubs Premium, zonal, con geo-replicación y clientes del protocolo Kafka. & Un espacio de nombres por ámbito, 2 en total; unidades y particiones según la memoria de capacidad. & CLD-01, con geo-replicación a CLD-02. Ámbito Retail y Emisor. & Conserva y permite releer el historial de hechos confirmados, con tres copias entre zonas y retención de hasta 90 días. Kafka en AKS exigiría operación propia y Service Bus no conserva historial. & SD-04, Tabla 4.6 y ADR-06 \\ \addlinespace
1.4 & Entrada de API & Azure API Management Premium clásico. & Una instancia por ámbito, 2 en total; 3 unidades por instancia, una por zona, 6 en total; un gateway adicional por instancia en CLD-02. & CLD-01 y gateway en CLD-02. Ámbito Retail y Emisor. & Único nivel con despliegue multirregión; los niveles v2 no están en Chile Central. Aplica identidad, cuotas, validación de contratos y versionado. & SD-04, Tabla 4.6 y ADR-06 \\ \addlinespace
1.5 & Datos transaccionales & Azure Database for PostgreSQL Flexible Server, nivel De uso general o Con optimización de memoria, con alta disponibilidad entre zonas. & Un servidor primario y uno en espera por ámbito, 4 nodos en total; una réplica entre regiones por ámbito aprobado; clase de cómputo según la memoria de capacidad. & CLD-01; réplica asíncrona en CLD-02. Ámbito Retail y Emisor. & Una base por autoridad del dato y un repositorio propio por servicio. La replicación síncrona entre zonas protege ante la falla de una zona. & SD-04, Tabla 4.6; nivel T-11 \\ \addlinespace
1.6 & Caché & Azure Managed Redis, servicio que sucede a Azure Cache for Redis. & Una instancia. & CLD-01. Ámbito Retail. & Guarda solo datos reconstruibles y no almacena autoridad de negocio. La sustitución del producto se explica bajo esta tabla. & SD-04, Tabla 4.6; producto T-11 \\ \addlinespace
1.7 & Almacenamiento analítico & Azure Data Lake Storage Gen2, con redundancia entre zonas. & Una cuenta por ámbito, 2 en total. & CLD-01; copia a CLD-02. Ámbito Retail y Emisor. & Separa la lectura analítica del almacenamiento transaccional, con control de acceso, linaje y exportación por entidad. & SD-04, Tabla 4.6 \\ \addlinespace
1.8 & Autoservicio de análisis & Power BI. & Un entorno por ámbito, 2 en total; licencias según usuarios habilitados. & CLD-01. Ámbito Retail y Emisor. & Espacios de lectura separados por entidad. El motor analítico de Azure se selecciona tras la prueba de permisos y linaje, porque Fabric completo no está disponible en Chile Central. & SD-04, Tabla 4.6 \\ \addlinespace
1.9 & Evidencia y expediente del Emisor & Azure Blob Storage con política de inmutabilidad. & Una cuenta. & CLD-01; copia a CLD-02 sujeta a la aprobación de residencia. Ámbito Emisor. & Sostiene el Servicio de evidencia financiera (F:C-03) con retención por el plazo del crédito más seis años. & SD-04, Tabla 4.10; producto T-11 \\ \addlinespace
1.10 & Identidad & Microsoft Entra ID, nivel P2, federado con el directorio del Cliente. & Identidades segregadas por ámbito; usuarios según el Caso (6.400 personas con acceso, 8.300 en período de mayor dotación). & Servicio global administrado desde CLD-01. Ámbito Retail y Emisor. & Identidad centralizada con autenticación multifactor, acceso condicional y elevación de privilegios con vigencia limitada. & SD-04, Tabla 4.6; nivel T-11 \\ \addlinespace
1.11 & Secretos y claves & Azure Key Vault, nivel Premium, con claves protegidas por módulo de seguridad de hardware. & Una bóveda por ambiente de cada ámbito. & CLD-01; claves copiadas a CLD-02. Ámbito Retail y Emisor. & Custodia de claves con rotación y separación de funciones. Las claves del Emisor residen en bóvedas propias, protegidas por hardware, sin equipo dedicado en sitio. & SD-04, Tabla 4.6; nivel T-11 \\ \addlinespace
1.12 & Telemetría & Azure Monitor con OpenTelemetry. & Una configuración por ámbito. & CLD-01. Ámbito Retail y Emisor. & Observabilidad conjunta de nube, centro de datos y tiendas. & SD-04, Tabla 4.6 \\ \addlinespace
1.13 & Registros & Áreas de trabajo de Log Analytics. & Una por ámbito, 2 en total. & CLD-01. Ámbito Retail y Emisor. & Registro inalterable con 12 meses en línea y 24 en archivo, separado por entidad. & SD-04, 4.3.1 \\ \addlinespace
1.14 & Borde público & Azure Front Door, nivel Premium, con firewall de aplicaciones web. & Un perfil. & Servicio global. Ámbito Compartido. & Entrada única del tráfico público con protección de aplicación y contra denegación de servicio. Redirige a CLD-02 en la conmutación. & SD-04, 4.2; nivel T-11 \\ \addlinespace
1.15 & Control de red & Azure Firewall, nivel Premium. & Uno en el concentrador de cada región activa. & CLD-01; se crea en CLD-02. Ámbito Compartido. & Bloquea todo acceso del radio Retail a los datos del Emisor y permite solo el contrato del Servicio de control de cruces (X-01). & SD-04, 4.2; nivel T-11 \\ \addlinespace
1.16 & Resolución de nombres & DNS privado de Azure. & Zonas privadas por radio. & CLD-01 y CLD-02. Ámbito Compartido. & Resolución interna sin exponer nombres ni direcciones a Internet. & SD-04, 4.2 \\ \addlinespace
1.17 & Acceso privado a datos & Puntos de conexión privados de Azure Private Link. & Uno por servicio de datos en la subred de datos de cada radio. & CLD-01 y CLD-02. Ámbito Retail, Emisor y compartido. & Ningún servicio de datos es alcanzable desde Internet. & SD-04, 4.2 \\ \addlinespace
1.18 & Red virtual & Redes virtuales en topología de concentrador y radios. & Un concentrador y 3 radios (Retail, Emisor y plataforma compartida), cada uno con subredes de ingreso, aplicación y datos. & CLD-01; se replica en CLD-02. Ámbito Retail, Emisor y compartido. & Hace efectiva la separación de redes entre las dos entidades. & SD-04, 4.2 \\ \addlinespace
1.19 & Conexión privada con DC-01 & Circuito y puerta de enlace de Azure ExpressRoute. & Un circuito y una puerta de enlace zonal. & Concentrador de CLD-01 y DC-01. Ámbito Compartido. & Convivencia con las plataformas conservadas sin pasar por Internet. & SD-04, 4.2 \\ \addlinespace
1.20 & Respaldo de la conexión y sitios & Puerta de enlace VPN de Azure. & Una puerta de enlace zonal. & Concentrador de CLD-01. Ámbito Compartido. & Respaldo cifrado del circuito privado y terminación de los túneles de tiendas, centros de distribución y oficinas. & SD-04, 4.2 \\ \addlinespace
1.21 & Respaldos & Bóveda de Azure Backup con copia inmutable. & Una bóveda por ámbito. & CLD-01; copia a CLD-02 configurada de forma explícita. Ámbito Retail y Emisor. & Chile Central no tiene región emparejada, de modo que la copia entre regiones no es automática. & SD-04, 4.2 y 4.3.1 \\ \addlinespace
1.22 & Registro de imágenes & Azure Container Registry, nivel Premium. & Un registro con replicación a CLD-02. & CLD-01 y CLD-02. Ámbito Compartido. & Guarda las imágenes firmadas que despliega el canal de entrega y las mantiene disponibles en la recuperación. & Formulario T-10; nivel T-11 \\ \addlinespace
1.23 & Cómputo de modelos predictivos & Carga de trabajo separada sobre AKS para entrenamiento e inferencia. & Un grupo de nodos dedicado; tamaño según la memoria de capacidad. & CLD-01. Ámbito Retail. & Los modelos de conteo y merma recomiendan y no escriben existencias, con trazabilidad y supervisión conforme a RT-18. & SD-04, 4.1.1 \\
\end{longtable}
}

La Tabla~\mbox{\ref{tab:t11-2}} confirma que cada ámbito dispone de entrada de API, intermediario de eventos, base de datos, almacenamiento analítico, claves y registros propios, y que la plataforma compartida solo aporta borde, red e integración. El único servicio exclusivo de un ámbito por decisión de diseño es la caché, que atiende solo a Retail.

La fila 1.6 sustituye el producto que nombra la Tabla 4.6 del Subdocumento 4. Microsoft dejó de admitir la creación de instancias de Azure Cache for Redis para clientes nuevos en abril de 2026 y fijó su retiro para el 30 de septiembre de 2028, dentro del período contractual (Microsoft, 2026a). Azure Managed Redis es el servicio que lo reemplaza y cumple la misma función. Su creación en Chile Central se comprueba junto con los demás servicios (Tabla~\mbox{\ref{tab:t11-14}}).

La cantidad de unidades de API Management se deriva de la sección 4.2 del Subdocumento 4, que asigna tres unidades a cada instancia para cubrir las tres zonas. Con dos instancias, una por ámbito, resultan seis unidades en la región primaria. Las demás cantidades de capacidad se obtienen del método de dimensionamiento de esa misma sección, que parte de las cifras del Caso 09 y cubre un crecimiento de tres veces la volumetría inicial en tres años y una carga de prueba de 1,5 veces el máximo previsto.

\subsection{T-11.2 Plataforma y software de la solución}\label{t-11-2}

La solución nueva se construye con una pila única para los servicios y otra para las interfaces. La Tabla~\mbox{\ref{tab:t11-3}} lista el software de base y el software propio, con la versión ofertada y el lugar donde se ejecuta.

{\footnotesize
\begin{longtable}[]{@{}
  >{\raggedright\arraybackslash}p{(\columnwidth - 10\tabcolsep) * \real{0.0500}}
  >{\raggedright\arraybackslash}p{(\columnwidth - 10\tabcolsep) * \real{0.1400}}
  >{\raggedright\arraybackslash}p{(\columnwidth - 10\tabcolsep) * \real{0.2200}}
  >{\raggedright\arraybackslash}p{(\columnwidth - 10\tabcolsep) * \real{0.1700}}
  >{\raggedright\arraybackslash}p{(\columnwidth - 10\tabcolsep) * \real{0.2900}}
  >{\raggedright\arraybackslash}p{(\columnwidth - 10\tabcolsep) * \real{0.1300}}@{}}
\caption{Plataforma y software de la solución.}\label{tab:t11-3}\\
\toprule\noalign{}
\textbf{N.°} & \textbf{Componente} & \textbf{Producto y versión ofertada} & \textbf{Ubicación y ámbito} & \textbf{Justificación} & \textbf{Origen} \\
\midrule\noalign{}
\endfirsthead
\toprule\noalign{}
\textbf{N.°} & \textbf{Componente} & \textbf{Producto y versión ofertada} & \textbf{Ubicación y ámbito} & \textbf{Justificación} & \textbf{Origen} \\
\midrule\noalign{}
\endhead
\bottomrule\noalign{}
\multicolumn{6}{@{}p{\columnwidth}@{}}{\rule{0pt}{3.2ex}\textbf{Fuente:} Elaboración propia a partir del Subdocumento 4 (Tablas 4.2, 4.6, 4.7, 4.10 y 4.11).} \\
\endlastfoot
2.1 & Lenguaje de los servicios & Java 25 LTS, en una distribución OpenJDK con soporte. & AKS en CLD-01 y CLD-02. Ámbito Retail y Emisor. & Versión de soporte prolongado vigente; una sola pila para APIs y lógica transaccional reduce la carga de operación y de competencias. & SD-04, Tabla 4.6 \\ \addlinespace
2.2 & Marco de los servicios & Spring Boot, línea 4, en su versión menor con soporte al inicio de la construcción. & AKS en CLD-01 y CLD-02. Ámbito Retail y Emisor. & Marco compatible con Java 25 para contratos síncronos, eventos y acceso a datos. & SD-04, Tabla 4.6 \\ \addlinespace
2.3 & Lenguaje y biblioteca de interfaces & TypeScript y React, en sus versiones estables vigentes al inicio de la construcción. & Navegador, portales y vistas del punto de venta. Ámbito Retail y Emisor. & Una sola base de componentes para portales internos, externos y caja. & SD-04, Tabla 4.6 \\ \addlinespace
2.4 & Orquestación de contenedores & Kubernetes, en una versión admitida por AKS. & CLD-01 y CLD-02. Ámbito Retail y Emisor. & Aislamiento por entidad, escalado y recuperación con la misma definición. & SD-04, Tabla 4.6 \\ \addlinespace
2.5 & Motor relacional & PostgreSQL, en la versión mayor más reciente que el servicio administrado ofrezca en la región. & CLD-01 y CLD-02. Ámbito Retail y Emisor. & Motor abierto, portable y con alta disponibilidad administrada. & SD-04, Tabla 4.6 \\ \addlinespace
2.6 & Protocolo de mensajería & Clientes del protocolo Kafka contra Event Hubs. & AKS y adaptadores. Ámbito Retail y Emisor. & El código de publicación y consumo no depende de una interfaz propia del proveedor. & SD-04, ADR-06 \\ \addlinespace
2.7 & Instrumentación & OpenTelemetry. & Todos los servicios y el nodo local. Ámbito Retail y Emisor. & Estándar abierto de métricas, registros y trazas, independiente del destino. & SD-04, Tabla 4.6 \\ \addlinespace
2.8 & Infraestructura como código & Terraform. & Repositorio y canal de entrega. Ámbito Compartido. & Toda la infraestructura de nube se define en código versionado, lo que permite crear CLD-02 sin derivar de la configuración de producción. & SD-04, Tabla 4.6 \\ \addlinespace
2.9 & Servicios de negocio & Trece servicios propios: R:M-01 a R:M-03, R:V-01 a R:V-04, R:CL-01, R:CL-02, F:C-01 a F:C-03 y X-01. & AKS en CLD-01. Ámbito Retail, Emisor y frontera. & Desarrollo propio que implementa la partición de dominios (ADR-05). & SD-04, Tabla 4.2 \\ \addlinespace
2.10 & Adaptadores & Componentes propios en contenedores. & CLD-01, o DC-01 y CD-01 cuando la interfaz lo exija. Ámbito Retail y Emisor. & Capa anticorrupción hacia cada plataforma conservada; la ruta se fija con el inventario de las catorce interfaces. & SD-04, Tablas 4.6 y 4.7 \\ \addlinespace
2.11 & Software del punto de venta & Aplicación propia de caja con modo local. & Terminales de caja de SIT-01. Ámbito Retail. & Sustituye el punto de venta de 2014 tienda por tienda (ADR-04). & SD-04, Tabla 4.11 \\ \addlinespace
2.12 & Software del nodo local & Caché de oferta (L-01), diario local (L-02) y sincronizador (L-03), desarrollo propio. & Nodos locales de SIT-01. Ámbito Retail. & Sostiene la venta y el cobro durante 24 horas sin enlace y concilia al reconectar. & SD-04, Tablas 4.6 y 4.10 \\ \addlinespace
2.13 & Registro de salida transaccional & Patrón implementado en cada servicio productor. & AKS en CLD-01. Ámbito Retail y Emisor. & Ningún hecho se publica sin haberse confirmado en la base de su autoridad. & SD-04, 4.1 \\
\end{longtable}
}

La Tabla~\mbox{\ref{tab:t11-3}} muestra que los componentes de fabricante se concentran en la plataforma, mientras que los servicios de negocio, los adaptadores, el punto de venta y el nodo local son desarrollo propio cuya propiedad se transfiere al Cliente. El motor y el empaquetado del nodo local, junto con su sistema operativo, se fijan con el inventario de hardware y el ensayo de continuidad de 24 horas (Tabla~\mbox{\ref{tab:t11-14}}), que es la comprobación que la Tabla 4.6 exige antes de comprometerlos.

\subsection{T-11.3 Seguridad, entrega y operación}\label{t-11-3}

La cadena de entrega, el análisis de seguridad y el monitoreo usan las herramientas que declara el Formulario T-10, a las que este formulario agrega el producto de correlación de eventos y el de detección en puntos finales. La Tabla~\mbox{\ref{tab:t11-4}} las reúne con su función y su condición de soporte.

{\footnotesize
\begin{longtable}[]{@{}
  >{\raggedright\arraybackslash}p{(\columnwidth - 12\tabcolsep) * \real{0.0400}}
  >{\raggedright\arraybackslash}p{(\columnwidth - 12\tabcolsep) * \real{0.1400}}
  >{\raggedright\arraybackslash}p{(\columnwidth - 12\tabcolsep) * \real{0.1800}}
  >{\raggedright\arraybackslash}p{(\columnwidth - 12\tabcolsep) * \real{0.1200}}
  >{\raggedright\arraybackslash}p{(\columnwidth - 12\tabcolsep) * \real{0.1300}}
  >{\raggedright\arraybackslash}p{(\columnwidth - 12\tabcolsep) * \real{0.2600}}
  >{\raggedright\arraybackslash}p{(\columnwidth - 12\tabcolsep) * \real{0.1300}}@{}}
\caption{Herramientas de seguridad, entrega y operación.}\label{tab:t11-4}\\
\toprule\noalign{}
\textbf{N.°} & \textbf{Función} & \textbf{Producto} & \textbf{Ubicación} & \textbf{Condición de soporte} & \textbf{Justificación} & \textbf{Origen} \\
\midrule\noalign{}
\endfirsthead
\toprule\noalign{}
\textbf{N.°} & \textbf{Función} & \textbf{Producto} & \textbf{Ubicación} & \textbf{Condición de soporte} & \textbf{Justificación} & \textbf{Origen} \\
\midrule\noalign{}
\endhead
\bottomrule\noalign{}
\multicolumn{7}{@{}p{\columnwidth}@{}}{\rule{0pt}{3.2ex}\textbf{Fuente:} Elaboración propia a partir del Formulario T-10 y del Subdocumento 6.} \\
\endlastfoot
3.1 & Repositorio y canal de entrega & GitHub Enterprise y GitHub Actions, con ejecutores segregados para el Emisor. & Servicio en la nube; ejecutores en CLD-01 & Suscripción del fabricante & Revisión por pares, firma y despliegue auditado. Se reutiliza el canal del Cliente si satisface revisión, firma y trazabilidad. & Formulario T-10 \\ \addlinespace
3.2 & Análisis estático y calidad & SonarQube. & CLD-01 & Suscripción del fabricante & Bloquea el avance de una versión que no cumple los umbrales de calidad. & Formulario T-10 \\ \addlinespace
3.3 & Composición, imágenes e inventario de componentes & Trivy y Syft. & Canal de entrega & Comunidad, código abierto & Detecta vulnerabilidades en dependencias, imágenes e infraestructura como código, y genera el inventario de componentes. & Formulario T-10 \\ \addlinespace
3.4 & Firma de artefactos & Cosign. & Canal de entrega & Comunidad, código abierto & Solo se despliegan imágenes firmadas. & Formulario T-10 \\ \addlinespace
3.5 & Análisis dinámico & OWASP ZAP. & Ambiente de QA & Comunidad, código abierto & Prueba la aplicación en ejecución antes de Preproducción. & Formulario T-10 \\ \addlinespace
3.6 & Pruebas de contrato, extremo a extremo y carga & Pact, Playwright y k6. & QA y Preproducción & Comunidad, código abierto & Verifican contratos entre servicios, recorridos completos y la carga de 1,5 veces el máximo. & Formulario T-10 \\ \addlinespace
3.7 & Gestión del proyecto y documentación & Jira, Microsoft SharePoint y Slack. & Servicio en la nube & Suscripción del fabricante & Sistema de información del proyecto, ya comprometido. & Formularios T-9 y T-10 \\ \addlinespace
3.8 & Correlación de eventos de seguridad & Microsoft Sentinel sobre Log Analytics. & CLD-01, un área de trabajo por ámbito & Servicio administrado & Centraliza eventos de nube, centro de datos y tiendas en la misma plataforma de registros, sin operar un producto adicional. & T-11 \\ \addlinespace
3.9 & Protección de cargas y de puntos finales & Microsoft Defender for Cloud y Microsoft Defender for Endpoint. & Nube, nodos locales y estaciones de operación & Servicio administrado & Detección y respuesta en nube y en sitio con un solo proveedor de identidad y de registros. & T-11 \\ \addlinespace
3.10 & Configuración y despliegue local & Herramienta de configuración declarativa del nodo local, seleccionada entre las candidatas del Formulario T-10. & Nodos locales de SIT-01 & Según la herramienta seleccionada & Mantiene los 44 nodos con la misma configuración y permite reconstruir uno desde el repositorio. & Formulario T-10 \\
\end{longtable}
}

La Tabla~\mbox{\ref{tab:t11-4}} cubre los cuatro análisis que la matriz de cumplimiento compromete para el flujo de integración continua, que son el estático, el dinámico, el de composición y el de imágenes. Las filas 3.8 y 3.9 siguen el criterio de ADR-07, que privilegia la continuidad de un solo proveedor para identidad, datos y operación. La herramienta de la fila 3.10 se compromete con el diseño del nodo local, como establece el Formulario T-10.

\subsection{T-11.4 Hardware que provee el proponente}\label{t-11-4}

El único hardware que provee Only Simple Solutions es el nodo local de tienda, según la Tabla 4.11 del Subdocumento 4. La Tabla~\mbox{\ref{tab:t11-5}} lo especifica con sus repuestos.

\begin{longtable}[]{@{}
  >{\raggedright\arraybackslash}p{(\columnwidth - 4\tabcolsep) * \real{0.2000}}
  >{\raggedright\arraybackslash}p{(\columnwidth - 4\tabcolsep) * \real{0.5000}}
  >{\raggedright\arraybackslash}p{(\columnwidth - 4\tabcolsep) * \real{0.3000}}@{}}
\caption{Hardware provisto por Only Simple Solutions.}\label{tab:t11-5}\\
\toprule\noalign{}
\textbf{Atributo} & \textbf{4.1 Nodo local de tienda} & \textbf{4.2 Repuestos de nodo local} \\
\midrule\noalign{}
\endfirsthead
\toprule\noalign{}
\textbf{Atributo} & \textbf{4.1 Nodo local de tienda} & \textbf{4.2 Repuestos de nodo local} \\
\midrule\noalign{}
\endhead
\bottomrule\noalign{}
\multicolumn{3}{@{}p{\columnwidth}@{}}{\rule{0pt}{3.2ex}\textbf{Fuente:} Elaboración propia a partir del Subdocumento 4 (Tabla 4.11 y sección 4.2.1), del Subdocumento 6 y del Caso 09 (capítulo 15, p. 34).} \\
\endlastfoot
Marca y modelo de referencia & Servidor compacto de borde, Dell Technologies u opción técnicamente equivalente. & El mismo modelo de la fila 4.1, con configuración precargada. \\ \addlinespace
Cantidad & 44, a razón de 2 por tienda en 22 tiendas. & 5 \\ \addlinespace
Características mínimas & Equipo nuevo con garantía de fábrica; dos discos de estado sólido en espejo (RAID 1); doble fuente de alimentación; administración remota fuera de banda; cifrado de disco; configuración endurecida según CIS; procesador, memoria y disco dimensionados para 24 horas de operación de la tienda de mayor venta. & Las mismas de la fila 4.1. \\ \addlinespace
Redundancia & Dos nodos por tienda, uno como réplica; cada nodo tolera la falla de un disco. & No aplica \\ \addlinespace
Ubicación & Gabinete de cada tienda de SIT-01 & Bodega en Chile, con despacho a cada tienda \\ \addlinespace
Soporte & Crítico, 24x7 en sitio con resolución en 4 horas & Reemplazo en 4 horas desde el diagnóstico \\ \addlinespace
Justificación & La tienda vende y cobra sin enlace durante 24 horas (ADR-04), y la falla de un nodo no detiene la caja. & Existencia de reemplazo del 10 \% del parque de componentes críticos. \\
\end{longtable}

Las cantidades de la Tabla~\mbox{\ref{tab:t11-5}} se derivan del Caso. Dos nodos en cada una de las 22 tiendas suman 44, y el 10 \% de ese parque equivale a 4,4 unidades, que se redondean a 5. El Caso proyecta hasta 25 tiendas, lo que elevaría el parque a 50 nodos y los repuestos se mantendrían en 5. Se eligió el espejo de dos discos porque tolera la pérdida de un disco, se reconstruye sin detener el nodo y basta para el volumen de un día de operación, mientras que la réplica en el segundo nodo protege ante la pérdida del equipo completo. El modelo de referencia y la capacidad de procesador, memoria y disco se fijan con el volumen de datos de una tienda sin enlace, que el Caso entrega al proponente como cifra por estimar, y se confirman en el ensayo de continuidad (Tabla~\mbox{\ref{tab:t11-14}}).

\subsection{T-11.5 Hardware y obras que adquiere el Cliente}\label{t-11-5}

El Caso establece que el hardware de terreno, el equipamiento de red y las obras los adquiere y ejecuta el Cliente, y que el proponente especifica qué comprar, cuánto y con qué características (Multitiendas Ancoa S.A., 2026b, cap. 11, p. 26). La Tabla~\mbox{\ref{tab:t11-6}} entrega esa especificación. Only Simple Solutions configura y certifica cada equipo antes de la ola de despliegue de su tienda.

{\footnotesize
\begin{longtable}[]{@{}
  >{\raggedright\arraybackslash}p{(\columnwidth - 12\tabcolsep) * \real{0.0400}}
  >{\raggedright\arraybackslash}p{(\columnwidth - 12\tabcolsep) * \real{0.1300}}
  >{\raggedright\arraybackslash}p{(\columnwidth - 12\tabcolsep) * \real{0.1400}}
  >{\raggedright\arraybackslash}p{(\columnwidth - 12\tabcolsep) * \real{0.1300}}
  >{\raggedright\arraybackslash}p{(\columnwidth - 12\tabcolsep) * \real{0.2600}}
  >{\raggedright\arraybackslash}p{(\columnwidth - 12\tabcolsep) * \real{0.1100}}
  >{\raggedright\arraybackslash}p{(\columnwidth - 12\tabcolsep) * \real{0.1900}}@{}}
\caption{Especificación de hardware y obras de cargo del Cliente.}\label{tab:t11-6}\\
\toprule\noalign{}
\textbf{N.°} & \textbf{Componente} & \textbf{Referencia de marca} & \textbf{Cantidad} & \textbf{Características mínimas} & \textbf{Ubicación} & \textbf{Justificación} \\
\midrule\noalign{}
\endfirsthead
\toprule\noalign{}
\textbf{N.°} & \textbf{Componente} & \textbf{Referencia de marca} & \textbf{Cantidad} & \textbf{Características mínimas} & \textbf{Ubicación} & \textbf{Justificación} \\
\midrule\noalign{}
\endhead
\bottomrule\noalign{}
\multicolumn{7}{@{}p{\columnwidth}@{}}{\rule{0pt}{3.2ex}\textbf{Fuente:} Elaboración propia a partir del Subdocumento 4 (Tabla 4.11 y Figura 4.17) y del Caso 09 (capítulo 6, p. 14; capítulo 11, p. 26; capítulo 14, pp. 31-32; capítulo 15, pp. 33-36).} \\
\endlastfoot
5.1 & Conmutador de tienda & Opción técnicamente equivalente entre fabricantes con soporte en Chile. & 13 como mínimo, uno por tienda sin segmentación. & Administrable, con una VLAN por función (cajas, administración, Emisor, video y Wi-Fi de clientes), alimentación por Ethernet para periféricos y doble fuente. & Las 13 tiendas de SIT-01 sin segmentación & El Caso exige separar esas redes donde hoy no lo están y acreditar la separación de la red del Emisor. \\ \addlinespace
5.2 & Router y firewall de tienda & Fortinet u opción técnicamente equivalente. & 13 como mínimo, uno por tienda sin segmentación. & Dos puertos de enlace externo con conmutación automática, túneles cifrados hacia la nube, túnel independiente para el mesón financiero y gestión centralizada. & Las 13 tiendas de SIT-01 sin segmentación & Enlaza la tienda con los radios de nube y mantiene el tráfico del Emisor fuera de la red de Retail. \\ \addlinespace
5.3 & Gabinete de tienda & Gabinete cerrado con control de acceso. & 22, uno por tienda. & Espacio para 2 nodos, conmutador, router y firewall, y unidad de alimentación; sensores de temperatura y apertura integrados a la observabilidad. & SIT-01 & El Caso exige un gabinete por tienda dimensionado para sostener la autonomía sin enlace. \\ \addlinespace
5.4 & Alimentación ininterrumpida del gabinete & Unidad en línea de doble conversión, con gestión por red. & 22, una por gabinete. & Potencia igual a la suma del consumo de los 2 nodos, el conmutador y el router y firewall, con el margen de crecimiento del gabinete; autonomía que cubra el paso al respaldo eléctrico de la tienda y el apagado ordenado de los nodos. & SIT-01 & Un corte de energía no debe corromper el diario local ni interrumpir la sincronización. \\ \addlinespace
5.5 & Terminales de caja & Se conservan los existentes. & 380 líneas de caja, 430 en la proyección a tres años. & Compatibilidad con el software del punto de venta nuevo, verificada en el prototipo. & SIT-01 & El punto de venta se sustituye por software, tienda por tienda, sin cambiar el parque de cajas. \\ \addlinespace
5.6 & Periféricos de caja & Se conservan los existentes. & Un conjunto por línea de caja, 380. & Impresora fiscal, lector de código y terminal de pago, con controladores verificados en el prototipo. & SIT-01 & La Tabla 4.11 los mantiene a cargo del Cliente. \\ \addlinespace
5.7 & Terminales móviles de sala & Se conservan los existentes. & 640, compartidas por 3.820 vendedores. & Sistema operativo vigente o una de sus dos versiones anteriores, gestión remota de flota y operación sin conexión para los perfiles internos. & SIT-01 & El Caso no pide adquirirlas y sí especificar la dotación; la dotación requerida se calcula con el uso por turno medido en las tiendas piloto. \\ \addlinespace
5.8 & Impresoras de etiqueta & Se conservan las existentes. & Al menos una por tienda, 22. & Compatibilidad con la confirmación de actualización de precio. & Bodega de cada tienda de SIT-01 & Sostienen el cambio de precio físico mientras no haya etiqueta electrónica, que el Caso excluye de la instalación. \\ \addlinespace
5.9 & Equipamiento de los centros de distribución & Se conserva el existente. & Según el inventario de CD-01 y CD-02. & Lectores y terminales compatibles con el WMS conservado; en CD-02, equipo de captura durable local. & CD-01 y CD-02 & CD-02 opera con enlace único y registra localmente durante un corte. \\ \addlinespace
5.10 & Obras de sitio & No aplica. & Según el inventario por tienda. & Canalizaciones, circuitos eléctricos independientes y cableado estructurado certificado para el gabinete. & SIT-01, CD-01 y CD-02 & El Caso asigna la ejecución al Cliente y la especificación al proponente. \\
\end{longtable}
}

La Tabla~\mbox{\ref{tab:t11-6}} separa lo que se compra de lo que se conserva. Solo las filas 5.1 a 5.4 y 5.10 implican adquisición u obra, y se concentran en las 13 tiendas que hoy carecen de segmentación y en los 22 gabinetes. Las nueve tiendas que ya tienen red segmentada conservan su conmutador y su router si el inventario confirma que cumplen las características mínimas. Las cantidades finales por tienda y los modelos de referencia se entregan en la especificación de hardware y dispositivos de terreno, que es el paquete 1.4.8 de la EDT (Tabla~\mbox{\ref{tab:t11-14}}). Los equipos críticos de esta tabla llevan soporte 24x7 con resolución en 4 horas y una existencia de reemplazo del 10 \% en sitio con configuración precargada.

\subsection{T-11.6 Conectividad}\label{t-11-6}

La conectividad une los sitios del Cliente con la nube por rutas privadas o cifradas. La Tabla~\mbox{\ref{tab:t11-7}} lista los enlaces por sitio y distingue los existentes de los que la propuesta requiere.

{\footnotesize
\begin{longtable}[]{@{}
  >{\raggedright\arraybackslash}p{(\columnwidth - 10\tabcolsep) * \real{0.0500}}
  >{\raggedright\arraybackslash}p{(\columnwidth - 10\tabcolsep) * \real{0.1500}}
  >{\raggedright\arraybackslash}p{(\columnwidth - 10\tabcolsep) * \real{0.2000}}
  >{\raggedright\arraybackslash}p{(\columnwidth - 10\tabcolsep) * \real{0.2500}}
  >{\raggedright\arraybackslash}p{(\columnwidth - 10\tabcolsep) * \real{0.1200}}
  >{\raggedright\arraybackslash}p{(\columnwidth - 10\tabcolsep) * \real{0.2300}}@{}}
\caption{Enlaces por sitio.}\label{tab:t11-7}\\
\toprule\noalign{}
\textbf{N.°} & \textbf{Sitio} & \textbf{Enlace existente} & \textbf{Enlace requerido por la solución} & \textbf{Cantidad nueva} & \textbf{Provee} \\
\midrule\noalign{}
\endfirsthead
\toprule\noalign{}
\textbf{N.°} & \textbf{Sitio} & \textbf{Enlace existente} & \textbf{Enlace requerido por la solución} & \textbf{Cantidad nueva} & \textbf{Provee} \\
\midrule\noalign{}
\endhead
\bottomrule\noalign{}
\multicolumn{6}{@{}p{\columnwidth}@{}}{\rule{0pt}{3.2ex}\textbf{Fuente:} Elaboración propia a partir del Subdocumento 4 (sección 4.2 y Tabla 4.9) y del Caso 09 (capítulo 6, p. 14; capítulo 15, p. 34).} \\
\endlastfoot
6.1 & 8 tiendas en centro comercial & Infraestructura del centro comercial y enlace propio de respaldo. & Túnel cifrado hacia la nube sobre ambos enlaces. & 0 & Cliente, enlaces; proponente, configuración \\ \addlinespace
6.2 & 6 tiendas en centro comercial & Infraestructura del centro comercial, sin respaldo propio. & Enlace de respaldo por un camino y un proveedor distintos. & 6 & Cliente \\ \addlinespace
6.3 & 8 tiendas a la calle & Enlace propio de un proveedor. & Enlace de respaldo por un camino y un proveedor distintos. & 8 & Cliente \\ \addlinespace
6.4 & CD-01 & Doble enlace con dos proveedores. & Túnel cifrado hacia la nube sobre ambos enlaces. & 0 & Cliente, enlaces; proponente, configuración \\ \addlinespace
6.5 & CD-02 & Enlace único. & Segundo enlace por un proveedor distinto. & 1 & Cliente \\ \addlinespace
6.6 & DC-01 & Sin conexión privada con la nube. & Circuito ExpressRoute y VPN de respaldo. & 1 circuito y 1 túnel & Proponente \\ \addlinespace
6.7 & DC-01 y DC-02 & Sin enlace acreditado entre ambos. & Enlace entre el centro de datos y su sala de respaldo. & 1 & Cliente \\ \addlinespace
6.8 & Oficinas & Acceso existente. & Túnel cifrado o acceso por red definida por software. & 0 & Cliente, enlaces; proponente, configuración \\
\end{longtable}
}

La Tabla~\mbox{\ref{tab:t11-7}} suma 14 enlaces de respaldo nuevos en tiendas. El Caso exige respaldo en las 14 tiendas de centro comercial y en Coyhaique, de modo que el mínimo exigible son 7, que corresponden a las 6 tiendas de centro comercial sin respaldo propio y a Coyhaique. Los otros 7 obedecen al supuesto SUP-28 del Subdocumento 2, según el cual ninguna tienda a la calle tiene respaldo. El levantamiento de enlaces por tienda confirma o reduce esa cifra. El ancho de banda de cada enlace, normal y en el máximo, sale del método de dimensionamiento de la sección 4.2.

\subsection{T-11.7 Centros de datos}\label{t-11-7}

La región primaria y la región de recuperación de Azure cumplen el papel de centro de datos principal y secundario de la solución nueva. La Tabla~\mbox{\ref{tab:t11-8}} compara sus especificaciones y agrega las instalaciones del Cliente, que permanecen en convivencia.

\begin{longtable}[]{@{}
  >{\raggedright\arraybackslash}p{(\columnwidth - 6\tabcolsep) * \real{0.1600}}
  >{\raggedright\arraybackslash}p{(\columnwidth - 6\tabcolsep) * \real{0.2700}}
  >{\raggedright\arraybackslash}p{(\columnwidth - 6\tabcolsep) * \real{0.2700}}
  >{\raggedright\arraybackslash}p{(\columnwidth - 6\tabcolsep) * \real{0.3000}}@{}}
\caption{Especificaciones de los centros de datos.}\label{tab:t11-8}\\
\toprule\noalign{}
\textbf{Atributo} & \textbf{Primario, CLD-01} & \textbf{Secundario, CLD-02} & \textbf{Instalaciones del Cliente, DC-01 y DC-02} \\
\midrule\noalign{}
\endfirsthead
\toprule\noalign{}
\textbf{Atributo} & \textbf{Primario, CLD-01} & \textbf{Secundario, CLD-02} & \textbf{Instalaciones del Cliente, DC-01 y DC-02} \\
\midrule\noalign{}
\endhead
\bottomrule\noalign{}
\multicolumn{4}{@{}p{\columnwidth}@{}}{\rule{0pt}{3.2ex}\textbf{Fuente:} Elaboración propia a partir del Subdocumento 4 (sección 4.3 y Tabla 4.8) y del Caso 09 (capítulo 6, p. 14).} \\
\endlastfoot
Identificación & Región Azure Chile Central (\texttt{chilecentral}). & Región Azure Brazil South (\texttt{brazilsouth}). & Centro de datos de casa matriz de 140 m² y sala de respaldo. \\ \addlinespace
Ubicación & Santiago, Chile. & São Paulo, Brasil. & Misma comuna ambos. \\ \addlinespace
Zonas de disponibilidad & Tres. & Tres. & No aplica. \\ \addlinespace
Función & Producción de los servicios nuevos. & Recuperación en modalidad activo-pasivo. & Alojamiento de las plataformas conservadas. \\ \addlinespace
Distancia y amenazas comunes & Países distintos, separados por 2.600 km aproximadamente. & Países distintos, separados por 2.600 km aproximadamente. & Comparten comuna, por lo que DC-02 no basta como única recuperación. \\ \addlinespace
Componentes & Filas 1.1 a 1.23 de la Tabla~\mbox{\ref{tab:t11-2}}. & Gateway de API Management, geo-réplica de Event Hubs, réplica de PostgreSQL, copia de almacenamiento, claves y respaldos, y AKS creado desde la definición de infraestructura. & Equipos del Cliente; el proponente aporta software de adaptadores y la conexión privada. \\ \addlinespace
Replicación & Síncrona entre zonas. & Asíncrona desde CLD-01, con retraso medido y alertado. & Enlace entre ambos recintos requerido (fila 6.7). \\ \addlinespace
Objetivos de recuperación & No aplica. & RTO no superior a 4 horas y RPO no superior a 15 minutos para servicios críticos. & No sostienen funciones nuevas como única ubicación. \\ \addlinespace
Ensayo & Falla de zona simulada. & Conmutación y retorno dos veces por año. & Conmutación de ExpressRoute a VPN. \\ \addlinespace
Energía, climatización y seguridad física & A cargo del proveedor de nube, con certificaciones entregadas cada año. & A cargo del proveedor de nube, con certificaciones entregadas cada año. & A cargo del Cliente; el plan de cierre de brecha se entrega en el hito H3. \\
\end{longtable}

La Tabla~\mbox{\ref{tab:t11-8}} muestra que ningún componente nuevo depende de un recinto del Cliente para recuperarse. El proponente no instala servidores nuevos en DC-01 ni en DC-02. Si el inventario de las catorce interfaces obliga a ejecutar un adaptador junto a una plataforma conservada, ese adaptador corre como contenedor sobre equipo del Cliente o sobre un nodo de las mismas características de la fila 4.1. La réplica y el ambiente de recuperación del Emisor en CLD-02 se activan con la aprobación de residencia de sus datos; mientras tanto, el Emisor se protege con la redundancia entre las tres zonas de CLD-01 y con respaldos inmutables en Chile.

\subsection{T-11.8 Plataformas conservadas}\label{t-11-8}

Las plataformas actuales del Cliente no forman parte de lo ofertado. La Tabla~\mbox{\ref{tab:t11-9}} las lista porque la solución se integra con ellas y porque su hardware y sus licencias no se atribuyen a los servicios nuevos.

\begin{longtable}[]{@{}
  >{\raggedright\arraybackslash}p{(\columnwidth - 6\tabcolsep) * \real{0.2500}}
  >{\raggedright\arraybackslash}p{(\columnwidth - 6\tabcolsep) * \real{0.2500}}
  >{\raggedright\arraybackslash}p{(\columnwidth - 6\tabcolsep) * \real{0.2000}}
  >{\raggedright\arraybackslash}p{(\columnwidth - 6\tabcolsep) * \real{0.3000}}@{}}
\caption{Plataformas del Cliente y su tratamiento.}\label{tab:t11-9}\\
\toprule\noalign{}
\textbf{Plataforma} & \textbf{Destino} & \textbf{Alojamiento supuesto} & \textbf{Aporte del proponente} \\
\midrule\noalign{}
\endfirsthead
\toprule\noalign{}
\textbf{Plataforma} & \textbf{Destino} & \textbf{Alojamiento supuesto} & \textbf{Aporte del proponente} \\
\midrule\noalign{}
\endhead
\bottomrule\noalign{}
\multicolumn{4}{@{}p{\columnwidth}@{}}{\rule{0pt}{3.2ex}\textbf{Fuente:} Elaboración propia a partir del Subdocumento 4 (Tablas 4.4 y 4.11).} \\
\endlastfoot
ERP y emisión de documentos tributarios & Se conserva como único emisor tributario. & DC-01 & Adaptador y modalidad fiscal de contingencia. \\ \addlinespace
WMS del centro de distribución principal & Se conserva; su extensión a Concepción se evalúa. & CD-01 & Adaptador y degradación por interfaz. \\ \addlinespace
Marketplace de 2022 & Se conserva. & Externo & Adaptador del Servicio de marketplace (R:V-04). \\ \addlinespace
Comercio electrónico de 2019 & Se integra durante la evaluación de Etapa 1. & Externo & Adaptador y pruebas de conservación o sustitución. \\ \addlinespace
Fidelización de 2017 & Se integra durante la evaluación de Etapa 1. & Externo & Adaptador y pruebas de conservación o sustitución. \\ \addlinespace
Núcleo Retail de 2009 & Se reemplaza por olas conciliadas. & DC-01 & Servicios de Retail de la fila 2.9. \\ \addlinespace
Plataforma de crédito de 2011 & Se reemplaza por olas de cartera y funciones. & DC-01 & Servicios del Emisor de la fila 2.9. \\ \addlinespace
Punto de venta de 2014 & Se sustituye tienda por tienda. & Terminales de SIT-01 & Filas 2.11 y 2.12. \\ \addlinespace
Motor de precios y promociones & No existe como sistema; es capacidad nueva. & No aplica & Servicio de oferta comercial (R:M-01). \\
\end{longtable}

La Tabla~\mbox{\ref{tab:t11-9}} delimita la oferta. El Cliente conserva las licencias y el hardware de las plataformas que se mantienen, y el proponente aporta solo los adaptadores y la integración. El alojamiento de cada plataforma se confirma con el inventario, y un cambio de ubicación modifica la ruta del adaptador respectivo sin alterar los servicios nuevos.

\subsection{T-11.9 Componentes de las innovaciones}\label{t-11-9}

Las cinco innovaciones del Subdocumento 13 agregan equipos y plataformas propios, que entran en producción en el mes 21, salvo el conteo predictivo, que opera desde el mes 16. La Tabla~\mbox{\ref{tab:t11-10}} los reúne para que la arquitectura, el plan y la oferta económica usen la misma lista.

{\footnotesize
\begin{longtable}[]{@{}
  >{\raggedright\arraybackslash}p{(\columnwidth - 10\tabcolsep) * \real{0.1500}}
  >{\raggedright\arraybackslash}p{(\columnwidth - 10\tabcolsep) * \real{0.1700}}
  >{\raggedright\arraybackslash}p{(\columnwidth - 10\tabcolsep) * \real{0.2000}}
  >{\raggedright\arraybackslash}p{(\columnwidth - 10\tabcolsep) * \real{0.1400}}
  >{\raggedright\arraybackslash}p{(\columnwidth - 10\tabcolsep) * \real{0.1000}}
  >{\raggedright\arraybackslash}p{(\columnwidth - 10\tabcolsep) * \real{0.2400}}@{}}
\caption{Componentes de las innovaciones.}\label{tab:t11-10}\\
\toprule\noalign{}
\textbf{Innovación} & \textbf{Componente} & \textbf{Producto de referencia} & \textbf{Ubicación} & \textbf{Cantidad} & \textbf{Justificación} \\
\midrule\noalign{}
\endfirsthead
\toprule\noalign{}
\textbf{Innovación} & \textbf{Componente} & \textbf{Producto de referencia} & \textbf{Ubicación} & \textbf{Cantidad} & \textbf{Justificación} \\
\midrule\noalign{}
\endhead
\bottomrule\noalign{}
\multicolumn{6}{@{}p{\columnwidth}@{}}{\rule{0pt}{3.2ex}\textbf{Fuente:} Elaboración propia a partir del Subdocumento 13.} \\
\endlastfoot
Casilleros de retiro & Gabinete electrónico exterior con controlador. & Cleveron 355 u opción técnicamente equivalente. & Tiendas definidas en el Subdocumento 13 & Según el Subdocumento 13 & Documentado para uso a la intemperie y con configuración modular para mercancía voluminosa. \\ \addlinespace
Conteo cíclico predictivo & Modelo supervisado de árboles con boosting. & scikit-learn, sobre la carga de la fila 1.23. & CLD-01 & Un servicio & No requiere hardware en tienda ni identificación por radiofrecuencia. \\ \addlinespace
Búsqueda vectorial multimodal & Modelo de representación de imagen y texto. & Jina CLIP v2, con licencia comercial. & CLD-01 & Un servicio & Admite consulta por imagen y por texto en español. \\ \addlinespace
Búsqueda vectorial multimodal & Índice vectorial. & Extensión pgvector sobre el PostgreSQL de Retail (fila 1.5). & CLD-01 & Una base & Reutiliza el motor transaccional y evita operar un motor adicional. \\ \addlinespace
Red de medios & Plataforma de publicidad en sitio. & Criteo u opción técnicamente equivalente, con adaptador propio. & Servicio externo, EXT-01 & Un servicio & Plataforma especializada frente a un desarrollo propio. \\ \addlinespace
Modo Tienda con fila virtual & Tótem de autoservicio. & Zebra KC50 de 22 pulgadas, con Ethernet. & Tiendas definidas en el Subdocumento 13 & Según el Subdocumento 13 & Montaje accesible y administración empresarial; la fila corre en el nodo local durante un corte. \\ \addlinespace
Modo Tienda con fila virtual & Impresora de turnos. & Epson TM-T20III Ethernet. & Junto a cada tótem & Según el Subdocumento 13 & Entrega el turno en papel cuando el flujo lo requiere. \\
\end{longtable}
}

La Tabla~\mbox{\ref{tab:t11-10}} agrega dos cargas a componentes ya listados. La fila virtual se ejecuta en el nodo local durante los cortes y entra en su dimensionamiento, y el índice vectorial usa el PostgreSQL de Retail, cuya admisión de la extensión se comprueba en la suscripción.

\subsection{T-11.10 Licenciamiento}\label{t-11-10}

Las licencias de software de terceros se constituyen a nombre del Cliente, son transferibles y rigen durante los 56 meses del contrato. La Tabla~\mbox{\ref{tab:t11-11}} declara la modalidad y la licencia de cada producto, incluido el software de código abierto y su compatibilidad con el uso previsto.

\begin{longtable}[]{@{}
  >{\raggedright\arraybackslash}p{(\columnwidth - 8\tabcolsep) * \real{0.2200}}
  >{\raggedright\arraybackslash}p{(\columnwidth - 8\tabcolsep) * \real{0.1700}}
  >{\raggedright\arraybackslash}p{(\columnwidth - 8\tabcolsep) * \real{0.1800}}
  >{\raggedright\arraybackslash}p{(\columnwidth - 8\tabcolsep) * \real{0.1100}}
  >{\raggedright\arraybackslash}p{(\columnwidth - 8\tabcolsep) * \real{0.3200}}@{}}
\caption{Licenciamiento por producto.}\label{tab:t11-11}\\
\toprule\noalign{}
\textbf{Producto} & \textbf{Modalidad} & \textbf{Licencia} & \textbf{Titular} & \textbf{Compatibilidad con el uso previsto} \\
\midrule\noalign{}
\endfirsthead
\toprule\noalign{}
\textbf{Producto} & \textbf{Modalidad} & \textbf{Licencia} & \textbf{Titular} & \textbf{Compatibilidad con el uso previsto} \\
\midrule\noalign{}
\endhead
\bottomrule\noalign{}
\multicolumn{5}{@{}p{\columnwidth}@{}}{\rule{0pt}{3.2ex}\textbf{Fuente:} Elaboración propia.} \\
\endlastfoot
Servicios de Azure de la Tabla~\mbox{\ref{tab:t11-2}} & Suscripción por consumo y por capacidad reservada. & Contrato de servicios en línea de Microsoft. & Cliente & Uso comercial sin restricción. \\ \addlinespace
Microsoft Entra ID P2, Power BI, Sentinel y Defender & Suscripción por usuario o por volumen. & Contrato de servicios en línea de Microsoft. & Cliente & Uso comercial sin restricción. \\ \addlinespace
GitHub Enterprise, SonarQube, Jira, SharePoint y Slack & Suscripción por usuario. & Licencia comercial del fabricante. & Cliente & Uso comercial sin restricción. \\ \addlinespace
Java 25, distribución OpenJDK & Código abierto con soporte del distribuidor. & GPL versión 2 con excepción Classpath. & No requiere titular & La excepción permite ejecutar aplicaciones propias sin obligación de liberarlas. \\ \addlinespace
Spring Boot, TypeScript, Kubernetes, OpenTelemetry, Trivy, Syft, Cosign, OWASP ZAP y Playwright & Código abierto. & Apache 2.0. & No requiere titular & Permisiva; sin obligación de liberar el código propio. \\ \addlinespace
React y Pact & Código abierto. & MIT. & No requiere titular & Permisiva. \\ \addlinespace
PostgreSQL y pgvector & Código abierto, dentro del servicio administrado. & Licencia PostgreSQL. & No requiere titular & Permisiva. \\ \addlinespace
scikit-learn & Código abierto. & BSD de tres cláusulas. & No requiere titular & Permisiva. \\ \addlinespace
Terraform & Código disponible. & Business Source License 1.1. & No requiere titular & Admite el uso interno para administrar infraestructura propia; no se redistribuye. \\ \addlinespace
k6 & Código abierto. & AGPL versión 3. & No requiere titular & Se usa como herramienta de prueba, sin modificarla ni ofrecerla como servicio, por lo que no impone obligaciones sobre la solución. \\ \addlinespace
Jina CLIP v2 & Licencia comercial del modelo. & Licencia comercial; la versión pública es no comercial. & Cliente & El suministro incluye la licencia comercial. \\ \addlinespace
Plataforma de publicidad y software de los casilleros & Suscripción o licencia del fabricante. & Licencia comercial del fabricante. & Cliente & Uso comercial sin restricción. \\ \addlinespace
Servicios de negocio, adaptadores, punto de venta y nodo local & Desarrollo propio. & Propiedad transferida al Cliente. & Cliente & No contiene componentes preexistentes del proponente. \\
\end{longtable}

La Tabla~\mbox{\ref{tab:t11-11}} muestra que ninguna licencia de código abierto impone la liberación del código de la solución. Las dos licencias que exigen atención son la del modelo multimodal, que se adquiere en su modalidad comercial, y la de la herramienta de carga, que se usa sin modificar. El registro de las licencias a nombre del Cliente es el paquete 1.4.7 de la EDT y se completa en el hito H3.

\subsection{T-11.11 Ciclo de vida y soporte}\label{t-11-11}

Todo componente ofertado tiene soporte vigente del fabricante o de su comunidad. La Tabla~\mbox{\ref{tab:t11-12}} declara, para cada componente principal, la versión ofertada, el fin de soporte publicado y el plan de actualización durante el contrato.

\begin{longtable}[]{@{}
  >{\raggedright\arraybackslash}p{(\columnwidth - 6\tabcolsep) * \real{0.2200}}
  >{\raggedright\arraybackslash}p{(\columnwidth - 6\tabcolsep) * \real{0.1700}}
  >{\raggedright\arraybackslash}p{(\columnwidth - 6\tabcolsep) * \real{0.2900}}
  >{\raggedright\arraybackslash}p{(\columnwidth - 6\tabcolsep) * \real{0.3200}}@{}}
\caption{Versión, fin de soporte y plan de actualización.}\label{tab:t11-12}\\
\toprule\noalign{}
\textbf{Componente} & \textbf{Versión ofertada} & \textbf{Fin de soporte publicado} & \textbf{Plan de actualización en el contrato} \\
\midrule\noalign{}
\endfirsthead
\toprule\noalign{}
\textbf{Componente} & \textbf{Versión ofertada} & \textbf{Fin de soporte publicado} & \textbf{Plan de actualización en el contrato} \\
\midrule\noalign{}
\endhead
\bottomrule\noalign{}
\multicolumn{4}{@{}p{\columnwidth}@{}}{\rule{0pt}{3.2ex}\textbf{Fuente:} Elaboración propia a partir de las referencias citadas, consultadas en octubre de 2026.} \\
\endlastfoot
Java & 25 LTS. & Septiembre de 2030 en soporte principal y septiembre de 2033 en soporte extendido (Oracle, 2026). & Migración a la siguiente versión de soporte prolongado durante la operación, antes de septiembre de 2030. \\ \addlinespace
Spring Boot & Línea 4, versión menor vigente. & Cada versión menor recibe soporte comunitario por unos 13 meses; el de la versión 4.0 termina en diciembre de 2026 (Spring, 2026). & Actualización a cada versión menor dentro de su ventana de soporte, al menos una vez por año. \\ \addlinespace
Kubernetes en AKS & Versión admitida por el servicio. & El servicio admite las tres versiones menores más recientes (Microsoft, 2026b). & Actualización programada al menos una vez por año, fuera de las ventanas de congelamiento del Caso. \\ \addlinespace
PostgreSQL & Versión mayor más reciente en la región. & El servicio sigue el calendario de la comunidad, que retira la versión 17 en noviembre de 2029 y la 18 en noviembre de 2030 (Microsoft, 2026c; PostgreSQL Global Development Group, 2026). & Una actualización de versión mayor durante la operación, ensayada en Preproducción. \\ \addlinespace
Azure Managed Redis & Versión administrada por el servicio. & Sin retiro anunciado; reemplaza a Azure Cache for Redis, que se retira el 30 de septiembre de 2028 (Microsoft, 2026a). & Actualización administrada por el proveedor. \\ \addlinespace
API Management Premium clásico, Event Hubs Premium y demás servicios administrados & Versión administrada por el servicio. & Sin retiro anunciado a la fecha de consulta. & Revisión trimestral de los avisos de retiro del proveedor, con migración planificada con un año de anticipación (riesgo R-07 del Formulario T-16). \\ \addlinespace
TypeScript y React & Versiones estables vigentes. & Sin calendario fijo; soporte sobre la versión mayor vigente. & Actualización anual de componentes de base. \\ \addlinespace
Terraform y herramientas de la cadena de entrega & Versiones estables vigentes. & Según cada proyecto. & Actualización anual, con las versiones fijadas en el repositorio. \\ \addlinespace
Nodo local de tienda & Equipo nuevo. & Garantía y soporte del fabricante por los 56 meses. & Reposición de cargo del proponente si el modelo deja de tener soporte antes del mes 56; borrado seguro certificado y disposición final con gestor autorizado al retiro. \\
\end{longtable}

La Tabla~\mbox{\ref{tab:t11-12}} identifica tres actualizaciones que ocurrirán con certeza dentro del contrato, que son las versiones menores de Spring Boot, las versiones de Kubernetes y una versión mayor de PostgreSQL, además de la migración de Java antes de 2030. Todas quedan a cargo del proponente como parte de la mantención, se ensayan en Preproducción y respetan las ventanas de congelamiento del Caso. El calendario de ciclo de vida se mantiene en la cuenta 1.15.4 de la EDT.

\subsection{T-11.12 Portabilidad}\label{t-11-12}

La estrategia de reversibilidad distingue los componentes que se trasladan a otro proveedor sin reescritura de los que son propios de Azure. La Tabla~\mbox{\ref{tab:t11-13}} clasifica cada grupo e indica en qué consiste su migración.

\begin{longtable}[]{@{}
  >{\raggedright\arraybackslash}p{(\columnwidth - 4\tabcolsep) * \real{0.3000}}
  >{\raggedright\arraybackslash}p{(\columnwidth - 4\tabcolsep) * \real{0.2000}}
  >{\raggedright\arraybackslash}p{(\columnwidth - 4\tabcolsep) * \real{0.5000}}@{}}
\caption{Portabilidad de los componentes.}\label{tab:t11-13}\\
\toprule\noalign{}
\textbf{Componente} & \textbf{Clasificación} & \textbf{Qué implica migrarlo} \\
\midrule\noalign{}
\endfirsthead
\toprule\noalign{}
\textbf{Componente} & \textbf{Clasificación} & \textbf{Qué implica migrarlo} \\
\midrule\noalign{}
\endhead
\bottomrule\noalign{}
\multicolumn{3}{@{}p{\columnwidth}@{}}{\rule{0pt}{3.2ex}\textbf{Fuente:} Elaboración propia a partir del Subdocumento 4 (Tabla 4.6).} \\
\endlastfoot
Servicios de negocio y adaptadores en contenedores & Portable & Desplegar las mismas imágenes en otro Kubernetes. \\ \addlinespace
PostgreSQL & Portable & Exportar e importar con herramientas estándar del motor. \\ \addlinespace
Clientes del protocolo Kafka & Portable & Cambiar el punto de conexión hacia un intermediario compatible con Kafka. \\ \addlinespace
Instrumentación OpenTelemetry & Portable & Cambiar el destino de exportación. \\ \addlinespace
Datos del lago analítico & Portable & Copiar archivos en formatos abiertos. \\ \addlinespace
Software del punto de venta y del nodo local & Portable & No depende de la nube elegida. \\ \addlinespace
Definición de infraestructura en Terraform & Portable en la herramienta & Reescribir los módulos para el proveedor de destino. \\ \addlinespace
API Management & Propio del proveedor & Traducir políticas, productos y suscripciones a otra puerta de enlace. \\ \addlinespace
Event Hubs, en su configuración de geo-replicación & Propio del proveedor & Reconfigurar la replicación en el intermediario de destino. \\ \addlinespace
Front Door y Azure Firewall & Propio del proveedor & Reescribir reglas de borde y de red. \\ \addlinespace
Key Vault & Propio del proveedor & Exportar las claves exportables y regenerar las que el hardware no libera. \\ \addlinespace
Entra ID & Propio del proveedor & Federar con otro proveedor de identidad mediante OIDC y SAML. \\ \addlinespace
Monitor, Log Analytics, Sentinel y Defender & Propio del proveedor & Exportar registros y reconstruir consultas y alertas. \\ \addlinespace
Power BI & Propio del proveedor & Reconstruir los informes sobre los mismos datos. \\
\end{longtable}

La Tabla~\mbox{\ref{tab:t11-13}} muestra que el código de negocio, los datos y el software de tienda son portables, y que la dependencia del proveedor se concentra en el borde, la identidad, las claves y la observabilidad. El esfuerzo de migración de cada componente propio del proveedor se estima en el Plan de Reversibilidad, y el riesgo de bloqueo se trata como R-08 en el Formulario T-16.

\subsection{T-11.13 Comprobaciones previas a la construcción}\label{t-11-13}

El Subdocumento 4 exige una comprobación de aceptación para cada tecnología antes de construir. La Tabla~\mbox{\ref{tab:t11-14}} reúne las comprobaciones que fijan los niveles, las cantidades y los modelos de este formulario, con el paquete de la EDT que las ejecuta y su hito.

\begin{longtable}[]{@{}
  >{\raggedright\arraybackslash}p{(\columnwidth - 6\tabcolsep) * \real{0.3000}}
  >{\raggedright\arraybackslash}p{(\columnwidth - 6\tabcolsep) * \real{0.3000}}
  >{\raggedright\arraybackslash}p{(\columnwidth - 6\tabcolsep) * \real{0.2400}}
  >{\raggedright\arraybackslash}p{(\columnwidth - 6\tabcolsep) * \real{0.1600}}@{}}
\caption{Comprobaciones que fijan cantidades, niveles y modelos.}\label{tab:t11-14}\\
\toprule\noalign{}
\textbf{Comprobación} & \textbf{Qué fija en este formulario} & \textbf{Paquete o instancia} & \textbf{Hito} \\
\midrule\noalign{}
\endfirsthead
\toprule\noalign{}
\textbf{Comprobación} & \textbf{Qué fija en este formulario} & \textbf{Paquete o instancia} & \textbf{Hito} \\
\midrule\noalign{}
\endhead
\bottomrule\noalign{}
\multicolumn{4}{@{}p{\columnwidth}@{}}{\rule{0pt}{3.2ex}\textbf{Fuente:} Elaboración propia a partir del Subdocumento 4 (Tablas 4.6 a 4.9) y del Formulario T-14.} \\
\endlastfoot
Creación efectiva de cada recurso en Chile Central y en Brazil South. & Niveles de las filas 1.2 a 1.22, incluida la admisión de pgvector. & Verificación del socio certificado de Microsoft en la suscripción del proyecto. & Antes de H3, mes 6 \\ \addlinespace
Memoria de capacidad. & Nodos de AKS, unidades y particiones de Event Hubs, clase de cómputo de PostgreSQL, tamaño de la caché y ancho de banda por enlace. & Dimensionamiento de la sección 4.2. & H2, mes 4 \\ \addlinespace
Inventario por tienda y levantamiento de enlaces. & Cantidades finales de las filas 5.1 a 5.4 y 5.10, y de los enlaces 6.2 y 6.3. & Paquete 1.4.8 de la EDT. & Meses 3 a 6 \\ \addlinespace
Inventario de las catorce interfaces y de la ubicación de cada plataforma. & Ubicación de los adaptadores de la fila 2.10 y necesidad del enlace 6.7. & Levantamiento de integración. & Antes de H3, mes 6 \\ \addlinespace
Prototipo del punto de venta. & Empaquetado de la fila 2.11 y compatibilidad de las filas 5.5 y 5.6. & Prototipo de caja. & Antes del piloto, mes 6 \\ \addlinespace
Ensayo de continuidad de 24 horas en una tienda piloto. & Modelo de referencia, capacidad, motor, empaquetado y sistema operativo del nodo local, y potencia de la fila 5.4. & Piloto de corte y reconexión. & Meses 6 y 7 \\ \addlinespace
Prueba de permisos y linaje. & Motor analítico de la fila 1.8. & Prueba de analítica por ámbito. & Antes de Etapa 1 \\ \addlinespace
Aprobación de residencia de los datos del Emisor. & Réplica y ambiente de recuperación del Emisor en CLD-02. & Decisión del Cliente con su asesoría jurídica. & Antes de H3, mes 6 \\ \addlinespace
Informe interno de 2024 del centro de datos. & Alcance del plan de cierre de brecha de DC-01. & Paquetes 1.4.9 a 1.4.13 de la EDT. & H3, mes 6 \\
\end{longtable}

La Tabla~\mbox{\ref{tab:t11-14}} concentra las comprobaciones entre los meses 3 y 7, antes de la primera ola de despliegue en tiendas. Con ellas, el inventario adj-sd-04\_s2\_inventario-hardware.xlsx se completa por tienda y este formulario se actualiza bajo control de cambios, sin modificar la selección de productos ni el reparto de responsabilidades de la Tabla 4.11.

\subsection{Referencias}\label{t11-referencias}

Las siguientes fuentes respaldan los datos del Caso y las fechas de soporte citados en este formulario.

\noindent\hangindent=1.5em\hangafter=1 Microsoft. (2026a). \emph{What's new in Azure Cache for Redis}. \url{https://learn.microsoft.com/azure/azure-cache-for-redis/cache-whats-new}

\noindent\hangindent=1.5em\hangafter=1 Microsoft. (2026b). \emph{Supported Kubernetes versions in Azure Kubernetes Service (AKS)}. \url{https://learn.microsoft.com/azure/aks/supported-kubernetes-versions}

\noindent\hangindent=1.5em\hangafter=1 Microsoft. (2026c). \emph{Azure Database for PostgreSQL versioning policy}. \url{https://learn.microsoft.com/azure/postgresql/configure-maintain/concepts-version-policy}

\noindent\hangindent=1.5em\hangafter=1 Multitiendas Ancoa S.A. (2026b). \emph{Bases técnicas. Caso 09: Cadena multitienda} (Licitación N.° TFEP-01/2026, versión 1.0).

\noindent\hangindent=1.5em\hangafter=1 Oracle. (2026). \emph{Java SE Support Roadmap}. \url{https://www.oracle.com/java/technologies/java-se-support-roadmap.html}

\noindent\hangindent=1.5em\hangafter=1 PostgreSQL Global Development Group. (2026). \emph{Versioning policy}. \url{https://www.postgresql.org/support/versioning/}

\noindent\hangindent=1.5em\hangafter=1 Spring. (2026). \emph{Spring Boot support}. \url{https://spring.io/projects/spring-boot#support}

\ossFinalPage

\end{document}
